\setcounter{section}{14}






  \zwischenueberschrift{Übungsaufgaben}

\inputaufgabe
{}
{





Es sei $K$ ein
\definitionsverweis {algebraisch abgeschlossener Körper}{}{}
und sei

\mavergleichskette
{\vergleichskette
{R
}
{ = }{ K[X]
}
{  }{
}
{    }{
}
{   }{
}
}
{}{}{}
der
\definitionsverweis {Polynomring}{}{}
über $K$. Zeige, dass jede
\definitionsverweis {algebraische Funktion}{}{}
$f$ auf einer offenen Menge

\mavergleichskettedisp
{\vergleichskette
{U
}
{ =} { D(F)
}
{ \subseteq} { K\!-\!\operatorname{Spek}\, \left(   R   \right)
}
{ =} { {\mathbb A}^{1}_{ K }
}
{ } {
}
}
{}{}{}
die Form

\mavergleichskette
{\vergleichskette
{f
}
{ = }{ G/H
}
{  }{
}
{    }{
}
{   }{
}
}
{}{}{}
mit nicht kürzbaren

\mavergleichskette
{\vergleichskette
{ G,H
}
{ \in }{ R
}
{  }{
}
{    }{
}
{   }{
}
}
{}{}{}
und mit

\mavergleichskette
{\vergleichskette
{ D(F)
}
{ \subseteq }{ D(H)
}
{  }{
}
{    }{
}
{   }{
}
}
{}{}{}
besitzt.





}
{} {}

\inputaufgabegibtloesung
{}
{





Es sei $K$ ein
\definitionsverweis {algebraisch abgeschlossener Körper}{}{}
und sei $R$ eine
$K$-\definitionsverweis {Algebra von endlichem Typ}{}{,}
die ein
\definitionsverweis {faktorieller}{}{}
\definitionsverweis {Integritätsbereich}{}{}
sei. Zeige, dass jede
\definitionsverweis {algebraische Funktion}{}{}
$f$ auf einer offenen Menge

\mavergleichskettedisp
{\vergleichskette
{U
}
{ \subseteq} { K\!-\!\operatorname{Spek}\, \left(   R   \right)
}
{ } {
}
{ } {
}
{ } {
}
}
{}{}{}
die Form

\mavergleichskette
{\vergleichskette
{f
}
{ = }{ G/H
}
{  }{
}
{    }{
}
{   }{
}
}
{}{}{}
mit nicht kürzbaren

\mavergleichskette
{\vergleichskette
{ G,H
}
{ \in }{ R
}
{  }{
}
{    }{
}
{   }{
}
}
{}{}{}
und mit

\mavergleichskette
{\vergleichskette
{ U
}
{ \subseteq }{ D(H)
}
{  }{
}
{    }{
}
{   }{
}
}
{}{}{}
besitzt.





}
{} {}

\inputaufgabe
{}
{





Ergänze den Beweis zu

Lemma 14.4
.





}
{} {}

\inputaufgabe
{}
{





Es sei

\mavergleichskette
{\vergleichskette
{ U
}
{ \subseteq }{  K\!-\!\operatorname{Spek}\, \left(  R  \right)
}
{  }{
}
{    }{
}
{   }{
}
}
{}{}{}
eine
\definitionsverweis {offene Teilmenge}{}{}
im
$K$-\definitionsverweis {Spektrum}{}{}
von einer $K$-Algebra $R$ über einem
\definitionsverweis {algebraisch abgeschlossenen Körper}{}{}
$K$ und sei

\mavergleichskette
{\vergleichskette
{ f
}
{ \in }{ \Gamma (U, {\mathcal O} )
}
{  }{
}
{    }{
}
{   }{
}
}
{}{}{}
eine
\definitionsverweis {algebraische Funktion}{}{.}
Zeige, dass $f$ eine
\definitionsverweis {stetige Abbildung}{}{}
nach $K$ definiert.





}
{} {}

\inputaufgabe
{}
{





Zeige, dass der Ring $\Gamma (U, {\mathcal O} )$
\definitionsverweis {reduziert}{}{}
ist.





}
{} {}

\inputaufgabe
{}
{





Es sei $K$ ein
\definitionsverweis {algebraisch abgeschlossener Körper}{}{.}
Wir betrachten den  Punkt

\mavergleichskettedisp
{\vergleichskette
{ P
}
{ =} { (0,1)
}
{ \in} { V \left(  X^3-Y^3+1  \right)
}
{ =} { C
}
{ \subseteq} { {\mathbb A}^{2}_{ K }
}
}
{}{}{.}
Es sei

\mavergleichskette
{\vergleichskette
{ U
}
{ \defeq }{ C \setminus \{P\}
}
{  }{
}
{    }{
}
{   }{
}
}
{}{}{.}
Beschreibe eine
\definitionsverweis {algebraische Funktion}{}{}
auf $U$, die sich nicht zu einer algebraischen Funktion auf ganz $C$ ausdehnen lässt.





}
{} {}

\inputaufgabegibtloesung
{}
{





Wir betrachten die Neilsche Parabel

\mavergleichskettedisp
{\vergleichskette
{ C
}
{ =} { V \left(  Y^2-X^3  \right)
}
{ \subseteq} { {\mathbb A}^{2}_{ K }
}
{ } {
}
{ } {
}
}
{}{}{}
und den Punkt

\mavergleichskettedisp
{\vergleichskette
{ P
}
{ =} { (1,1)
}
{ \in} { C
}
{ } {
}
{ } {
}
}
{}{}{.}
Finde eine
\definitionsverweis {algebraische Funktion}{}{,}
die auf
\mathl{C \setminus \{P\}}{} definiert ist, aber nicht auf ganz $C$. Tipp: Finde unterschiedliche Faktorzerlegungen von
\mathl{X^3-X^2}{.}





}
{} {}

\inputaufgabe
{}
{





Es sei $K$ ein
\definitionsverweis {algebraisch abgeschlossener Körper}{}{}
und sei $R$ eine
$K$-\definitionsverweis {Algebra von endlichem Typ}{}{,}
die ein
\definitionsverweis {Integritätsbereich}{}{}
sei. Es sei

\mavergleichskette
{\vergleichskette
{ U
}
{ = }{ D( {\mathfrak a} )
}
{ \subseteq }{ K\!-\!\operatorname{Spek}\, \left(   R   \right)
}
{    }{
}
{   }{
}
}
{}{}{}
eine
\definitionsverweis {offene Menge}{}{}
zu einem
\definitionsverweis {Ideal}{}{}

\mavergleichskette
{\vergleichskette
{ {\mathfrak a}
}
{ = }{ \left(  f_1 , \ldots , f_n  \right)
}
{  }{
}
{    }{
}
{   }{
}
}
{}{}{.}
Zeige, dass für den Ring der
\definitionsverweis {algebraischen Funktionen}{}{}
die Gleichheit

\mavergleichskettedisp
{\vergleichskette
{ \Gamma ( U , {\mathcal O} )
}
{ =} { \bigcap_{i = 1}^n R_{f_i }
}
{ } {
}
{ } {
}
{ } {
}
}
{}{}{}
gilt, wobei der Durchschnitt im
\definitionsverweis {Quotientenkörper}{}{}
von $R$ genommen wird.





}
{} {}

\inputaufgabe
{}
{





Es sei $R$ eine kommutative
$K$-\definitionsverweis {Algebra von endlichem Typ}{}{}
über einem
\definitionsverweis {algebraisch abgeschlossenen Körper}{}{.}
Es sei

\mavergleichskette
{\vergleichskette
{ U
}
{ \subseteq }{ K\!-\!\operatorname{Spek}\, \left(  R  \right)
}
{  }{
}
{    }{
}
{   }{
}
}
{}{}{}
eine
\definitionsverweis {offene Teilmenge}{}{}
und
\maabb {f} { U } { K
} {}
eine Funktion. Es sei

\mavergleichskette
{\vergleichskette
{ U
}
{ = }{ \bigcup_{i \in I} U_i
}
{  }{
}
{    }{
}
{   }{
}
}
{}{}{}
eine
\definitionsverweis {offene Überdeckung}{}{}
mit der Eigenschaft, dass die Einschränkungen

\mavergleichskette
{\vergleichskette
{ f_i
}
{ = }{ f \!\mid_{ U_i }
}
{  }{
}
{    }{
}
{   }{}
}
{}{}{}
\definitionsverweis {algebraische Funktionen}{}{}
sind. Zeige, dass dann $f$ selbst algebraisch ist.





}
{} {}

\inputaufgabe
{}
{





Es sei $K$ ein
\definitionsverweis {algebraisch abgeschlossener Körper}{}{}
und sei $R$ eine
$K$-\definitionsverweis {Algebra von endlichem Typ}{}{,}
die ein
\definitionsverweis {Integritätsbereich}{}{}
sei. Zeige, dass zu
\definitionsverweis {offenen Mengen}{}{}

\mavergleichskette
{\vergleichskette
{U
}
{ \subseteq }{V
}
{  }{
}
{    }{
}
{   }{
}
}
{}{}{}
die
\definitionsverweis {Restriktionsabbildung}{}{}
\maabbdisp {} { \Gamma ( V , {\mathcal O} ) } { \Gamma ( U , {\mathcal O} )
} {}
\definitionsverweis {injektiv}{}{}
ist.





}
{} {}

\inputaufgabe
{}
{





Betrachte

\mavergleichskette
{\vergleichskette
{V
}
{ = }{ V(XW-YZ)
}
{ \subseteq }{ { {\mathbb A}_{ K }^{  4   } }
}
{    }{
}
{   }{
}
}
{}{}{.}
Beschreibe eine
\definitionsverweis {offene Menge}{}{}

\mavergleichskette
{\vergleichskette
{ U
}
{ \subseteq }{ { {\mathbb A}_{ K }^{  4   } }
}
{  }{
}
{    }{
}
{   }{
}
}
{}{}{}
derart, dass der zu

\mavergleichskette
{\vergleichskette
{ U \cap V
}
{ \subseteq }{ U
}
{  }{
}
{    }{
}
{   }{
}
}
{}{}{}
gehörende
\definitionsverweis {Ringhomomorphismus}{}{}
\maabbdisp {} { \Gamma (U, {\mathcal O} ) } { \Gamma ( U \cap V, {\mathcal O} )
} {}
nicht surjektiv ist.





}
{} {}

In der folgenden Aufgabe wird das Konzept des
\definitionsverweis {Limes einer Abbildung}{}{}
verwendet. Dabei könnte

Aufgabe 1.9

hilfreich sein.

\inputaufgabe
{}
{





Wir betrachten die durch

\mavergleichskettedisp
{\vergleichskette
{ C
}
{ =} { V \left(  Y^2-X^2-X^3  \right)
}
{ \subseteq} { {\mathbb A}^{2}_{ {\mathbb C} }
}
{ } {
}
{ } {
}
}
{}{}{}
gegebene Kurve, den Punkt

\mavergleichskette
{\vergleichskette
{ P
}
{ = }{ (0,0)
}
{ \in }{ C
}
{    }{
}
{   }{
}
}
{}{}{}
und das offene Komplement

\mavergleichskette
{\vergleichskette
{ U
}
{ = }{ C \setminus \{P\}
}
{  }{
}
{    }{
}
{   }{
}
}
{}{}{.}
\aufzaehlungdrei{Zeige, dass
\mathl{{ \frac{   Y  }{  X  } }}{} eine
\definitionsverweis {algebraische Funktion}{}{}
auf $U$ ist, die nicht auf ganz $C$ algebraisch ausdehnbar ist.
}{Zeige, dass der
\definitionsverweis {Abbildungslimes}{}{}
zur Funktion
\maabbdisp {\varphi = { \frac{   Y  }{  X  } }} { U } { {\mathbb C}
} {}
in $P$ nicht existiert.
}{Zeige, dass es
\definitionsverweis {Folgen}{}{}

\mathl{\left(  w_n  \right)_{n \in \N }}{} und
\mathl{\left(  z_n  \right)_{n \in \N }}{} in $U$ gibt, die beide gegen $P$
\definitionsverweis {konvergieren}{}{,}
für die die
\definitionsverweis {Bildfolgen}{}{}
unter $\varphi$ jeweils konvergieren, aber gegen unterschiedliche Werte.
}





}
{} {}

Wir kommen zu einer Reihe von
\zusatzklammer {in vielen Gebieten der Mathematik} {} {}
wichtigen Begriffen, die wesentliche Eigenschaften der Strukturgarbe auf einem $K$-Spektrum prägnant zusammenfassen.







Es sei $X$ ein
\definitionsverweis {topologischer Raum}{}{.}
Unter einer
\definitionswort {Prägarbe}{}
${ \mathcal F  }$ auf $X$ versteht man eine Zuordnung, die jeder
\definitionsverweis {offenen Menge}{}{}

\mavergleichskette
{\vergleichskette
{U
}
{ \subseteq }{ X
}
{  }{
}
{    }{
}
{   }{
}
}
{}{}{}
eine Menge ${ \mathcal  F   }  \left(   U   \right)$ und zu je zwei offenen Mengen

\mavergleichskette
{\vergleichskette
{U
}
{ \subseteq }{V
}
{  }{
}
{    }{
}
{   }{
}
}
{}{}{}
eine
\definitionsverweis {Abbildung}{}{}
\maabbdisp {\rho_{V,U}} { { \mathcal  F   }  \left(   V   \right)   } { { \mathcal  F   }  \left(   U   \right)
} {}
zuordnet, wobei diese Zuordnung die beiden folgenden Bedingungen erfüllen muss.
\aufzaehlungzwei {Zu

\mavergleichskette
{\vergleichskette
{U
}
{ = }{ V
}
{  }{
}
{    }{
}
{   }{
}
}
{}{}{}
ist

\mavergleichskettedisp
{\vergleichskette
{ \rho_{U,U}
}
{ =} {
\operatorname{Id}_{  { \mathcal  F   }  \left(   U   \right)  }
}
{ } {
}
{ } {
}
{ } {
}
}
{}{}{.}
} {Zu offenen Mengen

\mavergleichskettedisp
{\vergleichskette
{U
}
{ \subseteq} { V
}
{ \subseteq} { W
}
{ } {
}
{ } {
}
}
{}{}{}
ist stets

\mavergleichskettedisp
{\vergleichskette
{ \rho_{W,U}
}
{ =} { \rho_{V,U} \circ  \rho_{W,V}
}
{ } {
}
{ } {
}
{ } {
}
}
{}{}{.}
}






Die Abbildungen $\rho_{V,U}$ heißen dabei \stichwort {Restriktionsabbildungen} {.}







Eine
\definitionsverweis {Prägarbe}{}{}
${ \mathcal  F   }$ auf einem
\definitionsverweis {topologischen Raum}{}{}
$X$ heißt
\definitionswort {Prägarbe von Gruppen}{,}
wenn zu jeder
\definitionsverweis {offenen Menge}{}{}

\mavergleichskette
{\vergleichskette
{U
}
{ \subseteq }{ X
}
{  }{
}
{    }{
}
{   }{
}
}
{}{}{}
die Menge ${ \mathcal  F   }  \left(   U   \right)$ eine
\definitionsverweis {Gruppe}{}{}
und zu jeder Inklusion

\mavergleichskette
{\vergleichskette
{U
}
{ \subseteq }{ V
}
{  }{
}
{    }{
}
{   }{
}
}
{}{}{}
die Restriktionsabbildung
\maabbdisp {\rho_{V,U}} { { \mathcal  F   }  \left(   V   \right) } { { \mathcal  F   }  \left(   U   \right)
} {}
ein
\definitionsverweis {Gruppenhomomorphismus}{}{}
ist.











Eine
\definitionsverweis {Prägarbe}{}{}
${ \mathcal  F   }$ auf einem
\definitionsverweis {topologischen Raum}{}{}
$X$ heißt
\definitionswort {Prägarbe von kommutativen Ringen}{,}
wenn zu jeder
\definitionsverweis {offenen Menge}{}{}

\mavergleichskette
{\vergleichskette
{U
}
{ \subseteq }{ X
}
{  }{
}
{    }{
}
{   }{
}
}
{}{}{}
die Menge ${ \mathcal  F   }  \left(   U   \right)$ ein
\definitionsverweis {kommutativer Ring}{}{}
und zu jeder Inklusion

\mavergleichskette
{\vergleichskette
{U
}
{ \subseteq }{ V
}
{  }{
}
{    }{
}
{   }{
}
}
{}{}{}
die Restriktionsabbildung
\maabbdisp {\rho_{V,U}} { { \mathcal  F   }  \left(   V   \right) } { { \mathcal  F   }  \left(   U   \right)
} {}
ein
\definitionsverweis {Ringhomomorphismus}{}{}
ist.







\inputaufgabe
{}
{





Zeige, dass die Zuordnung, die jeder offenen Menge

\mavergleichskette
{\vergleichskette
{U
}
{ \subseteq }{ K\!-\!\operatorname{Spek}\, \left(   R   \right)
}
{  }{
}
{    }{
}
{   }{
}
}
{}{}{}
den
\definitionsverweis {Ring der algebraischen Funktionen}{}{}

\mathl{\Gamma ( U , {\mathcal O} )}{} und zu jeder Inklusion

\mavergleichskette
{\vergleichskette
{U_1
}
{ \subseteq }{ U_2
}
{  }{
}
{    }{
}
{   }{
}
}
{}{}{}
die Restriktionsabbildung
\zusatzklammer {siehe

Lemma 14.6
} {} {}
\maabbdisp {} { \Gamma (U_2 , {\mathcal O} ) } { \Gamma (U_1 , {\mathcal O} )
} {}
zuordnet, eine
\definitionsverweis {Prägarbe}{}{}
von
$K$-\definitionsverweis {Algebren}{}{}
ist.





}
{} {}







Es sei $X$ ein
\definitionsverweis {topologischer Raum}{}{.}
Unter einer
\definitionswort {Garbe}{}
${ \mathcal F  }$ auf $X$ versteht man eine
\definitionsverweis {Prägarbe}{}{}
${ \mathcal F  }$ auf $X$, die die folgenden Eigenschaften erfüllt.
\aufzaehlungzwei {Zu jeder
\definitionsverweis {offenen Überdeckung}{}{}

\mavergleichskette
{\vergleichskette
{ U
}
{ = }{ \bigcup_{ i \in I} U_i
}
{  }{
}
{    }{
}
{   }{
}
}
{}{}{}
und Elementen

\mavergleichskette
{\vergleichskette
{s,t
}
{ \in }{ { \mathcal  F   }  \left(   U   \right)
}
{  }{
}
{    }{
}
{   }{
}
}
{}{}{}
mit

\mavergleichskette
{\vergleichskette
{ \rho_{U, U_i } (s )
}
{ = }{\rho_{U, U_i } (t )
}
{  }{
}
{    }{
}
{   }{
}
}
{}{}{}
für alle

\mavergleichskette
{\vergleichskette
{i
}
{ \in }{ I
}
{  }{
}
{    }{
}
{   }{
}
}
{}{}{}
gilt

\mavergleichskette
{\vergleichskette
{s
}
{ = }{t
}
{  }{
}
{    }{
}
{   }{
}
}
{}{}{.}
} {Zu jeder offenen Überdeckung

\mavergleichskette
{\vergleichskette
{ U
}
{ = }{ \bigcup_{ i \in I} U_i
}
{  }{
}
{    }{
}
{   }{
}
}
{}{}{}
und Elementen

\mavergleichskette
{\vergleichskette
{s_i
}
{ \in }{ { \mathcal  F   }  \left(   U_i    \right)
}
{  }{
}
{    }{
}
{   }{
}
}
{}{}{}
mit

\mavergleichskette
{\vergleichskette
{  \rho_{U_i, U_i \cap U_j } \left(  s_i  \right)
}
{ = }{  \rho_{U_j, U_i \cap U_j } \left(  s_j  \right)
}
{  }{
}
{    }{
}
{   }{
}
}
{}{}{}
für alle

\mavergleichskette
{\vergleichskette
{i,j
}
{ \in }{ I
}
{  }{
}
{    }{
}
{   }{
}
}
{}{}{}
gibt es ein

\mavergleichskette
{\vergleichskette
{s
}
{ \in }{ { \mathcal  F   }  \left(   U   \right)
}
{  }{
}
{    }{
}
{   }{
}
}
{}{}{}
mit

\mavergleichskette
{\vergleichskette
{s_i
}
{ = }{ \rho_{U, U_i } (s)
}
{  }{
}
{    }{
}
{   }{
}
}
{}{}{}
für alle

\mavergleichskette
{\vergleichskette
{i
}
{ \in }{ I
}
{  }{
}
{    }{
}
{   }{
}
}
{}{}{.}
}







\inputaufgabe
{}
{





Es seien
\mathkor {} {X} {und} {Y} {}
\definitionsverweis {topologische Räume}{}{.}
Zu jeder
\definitionsverweis {offenen Teilmenge}{}{}

\mavergleichskette
{\vergleichskette
{U
}
{ \subseteq }{ X
}
{  }{
}
{    }{
}
{   }{
}
}
{}{}{}
betrachten wir

\mavergleichskettedisp
{\vergleichskette
{ { \mathcal  C   }  \left(   U   \right)
}
{ =} { C^0 (U,Y)
}
{ =} { \left\{  \varphi:U \rightarrow Y  \mid  \varphi \text{ stetig}         \right\}
}
{ } {
}
{ } {
}
}
{}{}{.}
Zeige, dass dies eine
\definitionsverweis {Garbe}{}{}
auf $X$ ist.





}
{} {}

\inputaufgabe
{}
{





Es sei $X$ ein
\definitionsverweis {topologischer Raum}{}{.}
Zu jeder
\definitionsverweis {offenen Teilmenge}{}{}

\mavergleichskette
{\vergleichskette
{U
}
{ \subseteq }{ X
}
{  }{
}
{    }{
}
{   }{
}
}
{}{}{}
betrachten wir

\mavergleichskettedisp
{\vergleichskette
{ { \mathcal  C   }  \left(   U   \right)
}
{ =} { C^0 (U,\R)
}
{ =} { \left\{  \varphi:U \rightarrow \R  \mid  \varphi \text{ stetig}         \right\}
}
{ } {
}
{ } {
}
}
{}{}{.}
Zeige, dass dies eine
\definitionsverweis {Garbe}{}{}
von kommutativen
$\R$-\definitionsverweis {Algebren}{}{}
auf $X$ ist.





}
{} {}

\inputaufgabe
{}
{





Es sei $M$ eine
\definitionsverweis {differenzierbare Mannigfaltigkeit}{}{.}
Zu jeder offenen Teilmenge

\mavergleichskette
{\vergleichskette
{ U
}
{ \subseteq }{ M
}
{  }{
}
{    }{
}
{   }{
}
}
{}{}{}
betrachten wir die Menge
\mathl{C^1(U,\R)}{} der
\definitionsverweis {differenzierbaren Funktionen}{}{}
auf $U$. Es sei

\mavergleichskette
{\vergleichskette
{ M
}
{ = }{ \bigcup_{i \in I} U_i
}
{  }{
}
{    }{
}
{   }{
}
}
{}{}{}
eine
\definitionsverweis {offene Überdeckung}{}{.}
\aufzaehlungdreiabc{Zeige, dass zu

\mavergleichskette
{\vergleichskette
{ V
}
{ \subseteq }{ U
}
{  }{
}
{    }{
}
{   }{
}
}
{}{}{}
offen und

\mavergleichskette
{\vergleichskette
{ f
}
{ \in }{ C^1(U,\R)
}
{  }{
}
{    }{
}
{   }{
}
}
{}{}{}
auch die
\definitionsverweis {Einschränkung}{}{}
$f \vert_V$ zu
\mathl{C^1(V,\R)}{} gehört.
}{Es sei

\mavergleichskette
{\vergleichskette
{ f
}
{ \in }{ C^1(M,\R)
}
{  }{
}
{    }{
}
{   }{
}
}
{}{}{.}
Zeige, dass

\mavergleichskette
{\vergleichskette
{ f
}
{ = }{ 0
}
{  }{
}
{    }{
}
{   }{
}
}
{}{}{}
genau dann ist, wenn sämtliche Einschränkungen

\mavergleichskette
{\vergleichskette
{ f \vert_{U_i }
}
{ = }{ 0
}
{  }{
}
{    }{
}
{   }{
}
}
{}{}{}
sind.
}{Es sei eine Familie

\mavergleichskette
{\vergleichskette
{ f_i
}
{ \in }{ C^1(U_i,\R)
}
{  }{
}
{    }{
}
{   }{
}
}
{}{}{}
von Funktionen gegeben, die die \anfuehrung{Verträglichkeitsbedingung}{}

\mavergleichskettedisp
{\vergleichskette
{ f_i \vert_{U_i \cap U_j }
}
{ =} { f_j \vert_{U_i \cap U_j }
}
{ } {
}
{ } {
}
{ } {
}
}
{}{}{}
für alle $i,j$ erfüllen. Zeige, dass es ein

\mavergleichskette
{\vergleichskette
{ f
}
{ \in }{ C^1(M,\R)
}
{  }{
}
{    }{
}
{   }{
}
}
{}{}{}
gibt mit

\mavergleichskette
{\vergleichskette
{ f \vert_{U_i }
}
{ = }{ f_i
}
{  }{
}
{    }{
}
{   }{
}
}
{}{}{}
für alle $i$.
}





}
{} {}

\inputaufgabe
{}
{





Zeige, dass die Zuordnung, die jeder offenen Menge

\mavergleichskette
{\vergleichskette
{U
}
{ \subseteq }{ K\!-\!\operatorname{Spek}\, \left(   R   \right)
}
{  }{
}
{    }{
}
{   }{
}
}
{}{}{}
den
\definitionsverweis {Ring der algebraischen Funktionen}{}{}

\mathl{\Gamma ( U , {\mathcal O} )}{} und zu jeder Inklusion

\mavergleichskette
{\vergleichskette
{U_1
}
{ \subseteq }{ U_2
}
{  }{
}
{    }{
}
{   }{
}
}
{}{}{}
die Restriktionsabbildung
\zusatzklammer {siehe

Lemma 14.6
} {} {}
\maabbdisp {} { \Gamma (U_2 , {\mathcal O} ) } { \Gamma (U_1 , {\mathcal O} )
} {}
zuordnet, eine
\definitionsverweis {Garbe}{}{}
von
$K$-\definitionsverweis {Algebren}{}{}
ist.





}
{} {}

In den folgenden Aufgaben werden Ultrafilter und minimale Primideale besprochen. Wir geben die Definitionen.





Ein \definitionsverweis {Primideal}{}{}
${\mathfrak p}$
in einem \definitionsverweis {kommutativen Ring}{}{}
heißt \definitionswort {minimales Primideal}{}, wenn es kein Primideal ${\mathfrak q}$ mit

\mavergleichskette
{\vergleichskette
{ {\mathfrak q}
}
{ \subset }{ {\mathfrak p}
}
{  }{
}
{    }{
}
{   }{
}
}
{}{}{}
gibt.











Es sei $R$ ein
\definitionsverweis {kommutativer Ring}{}{.}
Ein
\definitionsverweis {multiplikatives System}{}{}

\mavergleichskette
{\vergleichskette
{ F
}
{ \subseteq }{ R
}
{  }{
}
{    }{
}
{   }{
}
}
{}{}{}
nennt man einen \definitionswort {Ultrafilter}{,} wenn

\mavergleichskette
{\vergleichskette
{ 0
}
{ \notin }{ F
}
{  }{
}
{    }{
}
{   }{
}
}
{}{}{}
ist und wenn $F$ maximal mit dieser Eigenschaft ist.







\inputaufgabe
{}
{





Es sei $R$ ein
\definitionsverweis {kommutativer Ring}{}{}
und sei

\mavergleichskette
{\vergleichskette
{F
}
{ \subseteq }{ R
}
{  }{
}
{    }{
}
{   }{
}
}
{}{}{}
ein
\definitionsverweis {multiplikatives System}{}{}
mit

\mavergleichskette
{\vergleichskette
{ 0
}
{ \notin }{F
}
{  }{
}
{    }{
}
{   }{
}
}
{}{}{.}
Zeige, dass $F$ genau dann ein
\definitionsverweis {Ultrafilter}{}{}
ist, wenn es zu jedem

\mavergleichskette
{\vergleichskette
{g
}
{ \in }{ R
}
{  }{
}
{    }{
}
{   }{
}
}
{}{}{,}

\mavergleichskette
{\vergleichskette
{g
}
{ \notin }{F
}
{  }{
}
{    }{
}
{   }{
}
}
{}{}{,}
ein

\mavergleichskette
{\vergleichskette
{f
}
{ \in }{ F
}
{  }{
}
{    }{
}
{   }{
}
}
{}{}{}
und eine natürliche Zahl $n$ mit

\mavergleichskette
{\vergleichskette
{ fg^n
}
{ = }{ 0
}
{  }{
}
{    }{
}
{   }{
}
}
{}{}{}
gibt.





}
{} {}

\inputaufgabe
{}
{





Es sei $R$ ein
\definitionsverweis {kommutativer Ring}{}{}
und sei

\mavergleichskette
{\vergleichskette
{ F
}
{ \subseteq }{ R
}
{  }{
}
{    }{
}
{   }{
}
}
{}{}{}
ein
\definitionsverweis {Ultrafilter}{}{.}
Zeige, dass das Komplement von $F$ ein
\definitionsverweis {minimales Primideal}{}{}
in $R$ ist.





}
{} {}

\inputaufgabe
{}
{





Es sei $R$ ein
\definitionsverweis {kommutativer Ring}{}{}
und sei $S$ ein
\definitionsverweis {multiplikatives System}{}{}
mit

\mavergleichskette
{\vergleichskette
{ 0
}
{ \notin }{ S
}
{  }{
}
{    }{
}
{   }{
}
}
{}{}{.}
Zeige, dass $S$ in einem
\definitionsverweis {Ultrafilter}{}{}
enthalten ist.





}
{\zusatzklammer {Man benutze das Lemma von Zorn} {} {.}} {}

\inputaufgabe
{}
{





Es sei $R$ ein kommutativer,
\definitionsverweis {reduzierter}{}{}
Ring. Zeige, dass jeder
\definitionsverweis {Nullteiler}{}{}
in einem
\definitionsverweis {minimalen Primideal}{}{}
enthalten ist.





}
{} {}

\inputaufgabe
{}
{





Es sei $K$ ein
\definitionsverweis {algebraisch abgeschlossener Körper}{}{}
und sei

\mavergleichskette
{\vergleichskette
{ {\mathfrak a}
}
{ \subseteq }{ K[X_1 , \ldots , X_n ]
}
{  }{
}
{    }{
}
{   }{
}
}
{}{}{}
ein
\definitionsverweis {Radikal}{}{}
mit dem zugehörigen
\definitionsverweis {Restklassenring}{}{}

\mavergleichskette
{\vergleichskette
{R
}
{ = }{ K[X_1 , \ldots , X_n] / {\mathfrak a}
}
{  }{
}
{    }{
}
{   }{
}
}
{}{}{,}
der der
\definitionsverweis {Koordinatenring}{}{}
zu

\mavergleichskette
{\vergleichskette
{V
}
{ = }{ V( {\mathfrak a} )
}
{  }{
}
{    }{
}
{   }{
}
}
{}{}{}
ist. Zeige, dass die
\definitionsverweis {irreduziblen Komponenten}{}{}
von $V$ den
\definitionsverweis {minimalen Primidealen}{}{}
von $R$ entsprechen.





}
{} {}

\inputaufgabe
{}
{





Es sei $K$ ein algebraisch abgeschlossener Körper und sei $R$ eine kommutative $K$-Algebra von endlichem Typ. Zeige, dass die
\definitionsverweis {minimalen Primideale}{}{}
von $R$ den irreduziblen Komponenten von
\mathl{K\!-\!\operatorname{Spek}\, \left(  R  \right)}{} entsprechen.





}
{} {}





  \zwischenueberschrift{Aufgaben zum Abgeben}

\inputaufgabe
{3}
{





Es sei $K$ ein
\definitionsverweis {algebraisch abgeschlossener Körper}{}{,}
wir betrachten die affine Ebene ${\mathbb A}^{2}_{ K }$. Es sei

\mavergleichskette
{\vergleichskette
{ P
}
{ \in }{ {\mathbb A}^{2}_{ K }
}
{  }{
}
{    }{
}
{   }{
}
}
{}{}{}
ein Punkt und

\mavergleichskette
{\vergleichskette
{ U
}
{ = }{ {\mathbb A}^{2}_{ K }  \setminus \{P\}
}
{  }{
}
{    }{
}
{   }{
}
}
{}{}{}
das offene Komplement davon. Zeige

\mavergleichskettedisp
{\vergleichskette
{ \Gamma (U, {\mathcal O} )
}
{ =} { K[X,Y]
}
{ } {
}
{ } {
}
{ } {
}
}
{}{}{.}





}
{\zusatzklammer {Dies besagt, dass eine außerhalb eines Punktes der Ebene definierte algebraische Funktion sich in den Punkt fortsetzen lässt.} {} {.}} {}

\inputaufgabe
{5 (1+2+2)}
{





Wir betrachten die Neilsche Parabel

\mavergleichskettedisp
{\vergleichskette
{ C
}
{ =} { V \left(  X^2-Y^3  \right)
}
{ \subseteq} { {\mathbb A}^{2}_{ {\mathbb C} }
}
{ } {
}
{ } {
}
}
{}{}{.}
\aufzaehlungdreiabc{Zeige, dass auf $C$ die Gleichheit

\mavergleichskette
{\vergleichskette
{ D(X)
}
{ = }{ D(Y)
}
{  }{
}
{    }{
}
{   }{
}
}
{}{}{}
gilt.
}{Zeige, dass auf

\mavergleichskette
{\vergleichskette
{ U
}
{ = }{ D(Y)
}
{ \subseteq }{ C
}
{    }{
}
{   }{
}
}
{}{}{}
durch
\mathl{{ \frac{   X  }{  Y  } }}{} eine
\definitionsverweis {algebraische Funktion}{}{}
definiert ist, die sich nicht auf $C$ als algebraische Funktion ausdehnen lässt.
}{Zeige, dass die stetige Funktion
\maabbdisp {{ \frac{   X  }{  Y  } }} { D(Y) } { {\mathbb C}
} {}
eine
\definitionsverweis {stetige Fortsetzung}{}{}
auf ganz $C$ besitzt.
}





}
{} {}

\inputaufgabe
{4}
{





Es sei $R$ eine
\definitionsverweis {integre}{}{}
$K$-\definitionsverweis {Algebra}{}{}
$R$ von
\definitionsverweis {endlichem Typ}{}{}
über einem
\definitionsverweis {algebraisch abgeschlossenen Körper}{}{}
$K$. Es sei

\mavergleichskette
{\vergleichskette
{ q
}
{ \in }{ Q
}
{ = }{ Q(R)
}
{    }{
}
{   }{
}
}
{}{}{}
ein Element im
\definitionsverweis {Quotientenkörper}{}{}
von $R$. Zeige, dass

\mavergleichskettedisp
{\vergleichskette
{ {\mathfrak a}
}
{ =} { \left\{ f \in R  \mid  \text{es gibt } n \in \N \text{ mit } f^nq \in R        \right\}
}
{ } {
}
{ } {
}
{ } {
}
}
{}{}{}
ein
\definitionsverweis {Ideal}{}{}
in $R$ ist. Zeige ferner, dass

\mavergleichskette
{\vergleichskette
{ D({\mathfrak a})
}
{ \subseteq }{ K\!-\!\operatorname{Spek}\, \left(  R  \right)
}
{  }{
}
{    }{
}
{   }{
}
}
{}{}{}
der
\zusatzklammer {maximale} {} {}
Definitionsbereich der
\definitionsverweis {algebraischen Funktion}{}{}
$q$ ist.





}
{} {}

\inputaufgabe
{4}
{





Es sei $R$ ein
\definitionsverweis {kommutativer Ring}{}{}
und seien

\mavergleichskette
{\vergleichskette
{ f_1 , \ldots , f_n
}
{ \in }{ R
}
{  }{
}
{    }{
}
{   }{
}
}
{}{}{}
Elemente, die das
\definitionsverweis {Einheitsideal}{}{}
erzeugen. Es sei vorausgesetzt, dass die Nenneraufnahmen $R_{f_i }$ für

\mavergleichskette
{\vergleichskette
{ i
}
{ = }{ 1 , \ldots ,  n
}
{  }{
}
{    }{
}
{   }{
}
}
{}{}{}
\definitionsverweis {noethersch}{}{}
sind. Zeige, dass dann auch $R$ noethersch ist.





}
{} {}
